\section{\label{III_3_A}Statut des données : données automatiques et données manuelles/métiers}

L'introduction de processus fondés sur l'intelligence artificielle au sein d'un logiciel de gestion des collections modifie la nature et le statut des données produites sur les documents. Comme le souligne Jonas Arnold dans le \textit{Livre blanc sur l'apprentissage automatique dans les archives}, les données automatiquement produites possèdent une nature "fluide" puisqu'elles sont amenées à évoluer au rythme des améliorations des outils, internes ou externes \footcite{zotero-item-584}. 

(transition avec sous partie suivante : nécessité de généraliser le processus pour que les données ne bougent pas trop) (transition ausis avec refonte de S-Movies : le contrôle qualité exige une meilleure navigation de l'interface te une vue d'ensemble)

La nouvelle nature des données implique la mise en place d'une intelligence hybride qui impliquerait par exemple que l'utilisateur métier supervise la cohérence entre les entités du modèle. Ainsi, tandis que le processus automatique nourrit des entités inférieures, celles du fragment et de l'item, il reviendrait à l'utilisateur de garantir les liens avec les entités supérieures, celle de l'Oeuvre, et la cohérence du reversement des données au sein du modèle dans son ensemble. Le reversement des données automatique en \textit{backend} prévu au sein de la seconde partie est un revesrement théorique, qu'il revient à l'utilisateur de vérifier la cohérence. 
Afin d'assurer la vérification de la cohérence du modèle de données, deux dispositifs pourront être mis en place au sein du logiciel. 
D'une part, afin de pouvoir garantir l'homogénéité de l'indexation de ces données, le logiciel doit prévoir leur traçabilité via une "historisation ou un contrôle des versions" \footcite{zotero-item-584}. 
Skinsoft prévoit des champs de métadonnées de gestion au sein du logiciel contenant par exemple la date de modification de la notice, le dernier utilisateur à l'avoir modifié. Les métadonnées peuvent être définies comme de l'information structurée décrivant des données \footcite{colavizza2026}. Le même type de métadonnées pourra être conservé au sein de champs dédiés aux métadonnées de processus, telles que les outils utilisés, les paramètres appliqués, l'utilisateur ayant validé l'indexation, la date. Ces métadonnées apparaitraient dans une notice de Fragment par exemple. 
La traçablité des métadonnées produites permettrait ainsi à l'archiviste non seulement de suivre les évolutions et de veiller à la conformité de leur structure, mais également aux données d'être réemployées à des fins de réentraînement des machines. Ainsi, si le processus est modifié par l'emploi d'un nouvel outil par exemple, les métadonnées de processus permettront une réadaptation efficace sans modifier l'interface utilisateur encore davantage. 
D'autre part, pour une question de transparence et de visibilité au sein de l'interface, l'origine des données doit être excplicitement signalée. Si des métadonnées de processus sont intégrées au sein des notices, les données issues de processus automatiques doivent être clairement identifiée d'emblée au sein des grilles de résultats de recherche, par exemple. Cette identification peut prendre la forme d'un code couleur ou de tags de provenance apposés sur les données. Ce système permettrait à l'utilisateur de vérifier immédiatement si les données sont associées au bonne entités du modèle FRBR. 
Ainsi, la mise en place d'une intelligence hybride consiste à permettre la production automatique de données patrimoniales, à la condition que l'utilisateur garantisse leur cohérence et leur validation au sein du modèle de données. 
afin de gérer les erreurs potentielles et les bugs : système de tickets et de tests pour régler ensuite les régressions de l'application. 